\caption{%
  \textbf{\methodname{}: Retrieval-Augmented Diffusion for Data-Efficient Reasoning.}
  The paper section presents three core results.
  (1)~\methodname{} outperforms Baseline, Method-A, and Method-B on five reasoning
  benchmarks (MMLU, HumanEval, GSM8K, MATH, GPQA), with the largest gain of
  $+4.8$\,pp on MATH (Figure~1).
  (2)~\methodname{} is data-efficient, matching the next-best method's final
  validation loss in $60\%$ of the training steps ($\Delta{=}0.155$ at convergence;
  Figure~2).
  (3)~\methodname{} follows a Chinchilla-style scaling law ($y = -0.0449x + 0.621$,
  $R^{2}=0.877$), supporting the proposed power-law relationship between model size
  and loss (Figure~3).
  The architecture (Figure~4) consists of an Attention+Memory module, a Skill Router,
  and a residual-connected decoder; ablations confirm all four components are necessary.
  Data are synthetic but realistic (see \texttt{captions.json} and \texttt{claim.json}).%
}
